\section{Does an Ethical Criterion Generalize Beyond Its Training Distribution?}
\label{sec:11.3}

\textbf{Unresolved issue.}
A model may satisfy a fairness or refusal criterion during training without learning the principle the criterion was intended to represent. It remains unresolved whether an ethics-embedding method preserves that criterion under distribution shift or merely exploits correlations present in its training environment \cite{krakovna2020specification,langosco2022goal}.

\textbf{Test.}
Train matched models to satisfy the same ethical criterion on one distribution. Evaluate them on independently constructed distributions that vary one previously fixed feature at a time. Measure task capability and compliance with the ethical criterion separately, so that goal misgeneralization can be distinguished from ordinary capability loss.

\textbf{Evidence against the proposal.}
If task capability remains intact while the ethical criterion fails under modest shifts, the tested embedding method cannot support the proposed floor. Conversely, if ethical performance consistently degrades only when task capability does, the claim that ethics embedding creates a distinct generalization failure is not supported.

\textbf{Required access or authority.}
The test requires access to the training distribution, evaluation outputs, and enough model access to measure capability and criterion compliance independently. The shifted evaluation sets must be designed by researchers who did not develop the embedding method, and publication of all prespecified results must be guaranteed.
